%!TEX root = ../thesis.tex
%*******************************************************************************
%*********************************** First Chapter *****************************
%*******************************************************************************

\chapter{Introduction}  %Title of the First Chapter

\ifpdf
    \graphicspath{{Chapter1/Figs/Raster/}{Chapter1/Figs/PDF/}{Chapter1/Figs/}}
\else
    \graphicspath{{Chapter1/Figs/Vector/}{Chapter1/Figs/}}
\fi


\vspace{-1cm}

% % RL with MDP
% Reinforcement learning is machine learning algorithms that train an agent to behave optimally in a certain environment.
% In this thesis, we model the environment as a finite Markov decision process.
% The goal of the agent is then to learn the actions in each state that yield the highest total discounted rewards.

Reinforcement learning is a machine learning framework that trains an agent to behave optimally in a certain environment. More specifically, the agent learns to perform a task by choosing a sequence of actions that maximises the total reward it receives.

\vspace{-0.1cm}


% model-free PI
% In many real-world decision-making tasks, the parameters of the underlying MDP are unknown.
% The agent can then rely on the policy iteration (PI) framework where the agent repeatedly performs that task, in part or in full, and keeps track of the rewards it receives along the way. Doing so, it builds up experience about which action, or sequence of actions, yields higher rewards than others.
% It then uses this experience to progressively identify the optimal way to perform the task \citep{Silver2021RewardEnough}.

This framework has inspired many powerful algorithms over the last decades.
Among these, Monte Carlo (MC) methods stand out as they can train the agent without
% prior knowledge of the environment dynamics. 
knowing a model of the environment. 
Specifically, by repeatedly simulating the task, the agent accumulates experience and progressively identifies the optimal behaviour.
% simulating the task several times, the agent builds up experience which it then uses to progressively identify the optimal behaviour. 
% Thus these algorithms do not require prior knowledge of the environment dynamics. 
A striking application of MC methods is MuZero \citep{Schrittwieser2020MasteringModel}. This model achieved state-of-the-art performance in chess, Go, shogi and 57 Atari games using self-play without knowing the rules of each game.

\vspace{-0.1cm}


Despite these impressive results, MC methods in reinforcement learning lack a rigorous theoretical foundation. In particular, 
% the necessary and sufficient conditions for the global convergence of 
global convergence guarantees
for algorithms like MuZero and the foundational Monte Carlo optimistic policy iteration (MC-O-PI) are still unknown,
even when no function approximation is involved.
% This limitation hinders their widespread adoption in real-world applications.
% This is, according to renowned experts in the field,``one of the most fundamental open theoretical questions in reinforcement learning'' \citet{Sutton2018ReinforcementIntroduction}.

\vspace{-0.1cm}


% goal
Therefore, we investigate the convergence properties of MC-O-PI to address fundamental open questions and establish implementation guidelines.
We propose a new approach that studies the sequence of policies generated by the algorithm.
% averages out noise, which simplifies analysis and simulation.
This approach allows us to extend all existing convergence proofs, construct new counterexamples and conjecture the most general conditions for convergence to date.


This chapter first introduces finite Markov decision processes, the formal framework for agent-environment interactions.
It then presents the policy iteration framework, which underpins MC-O-PI,
before outlining the contributions of this thesis.

% \subsubsection{Notation}

\vspace{-0.1cm}


Unless stated otherwise,
% For notation,
operations and relations between functions are elementwise.
For instance, given functions $f, g : \xset \to \yset$
% for every operator $\star$ and relation $\bowtie$
we write
\begin{align*}
    &(f + g)(x) = f(x) + g(x), \ \forall \, x \in \xset ,
    \\
    % &(f \cdot g)(x) = f(x) \cdot g(x), \forall x \in \xset ,
    % \\
    &f \leq g \iff f(x) \leq g(x), \ \forall \, x \in \xset .
\end{align*}
Furthermore, every norm is computed as $\| f \| \coloneqq \sup_{x \in \xset} | f(x) | $.
% also convergence is pointwise
% \begin{align}
%     f_k \to f \implies \forall x : f_k(x) \to f(x)
% \end{align}
% denot sequence $a_0$, $a_1$, $a_2$, etc., by $(a_k)_{k \geq 0}$
% sometimes remove the subscript to alleviate notation

% Furthermore, if $\xset = \yset$, we refer to time-invariant dynamical systems as $x_+ = f(x)$, and time-variant ones as $x_\kpo = f(x_k)$.


% except stated otherwise, the norm $\|\cdot\|$ is always the infinity norm, so the elementwise maximum of its argument

% define constant functions $\onefun$ and $\zerofun$


\vspace{-0.2cm}

%%%%%%%%%%%%%%%%%%%%%%%%%%%%%%%%%%%%%%%%%
%%%%%%%%%%%%%%%%%%%%%%%%%%%%%%%%%%%%%%%%%
\section{Finite Markov decision processes}

\vspace{-0.2cm}

% \textcolor{red}{don't confuse random variable and realisation of random variable}


A finite Markov decision process (MDP) models decision-making using states, actions and rewards.
An agent interacts with its environment by selecting actions to transition between states while receiving rewards.
% that quantify its performance.
We now formally define the environment, the agent and their key properties.

\vspace{-0.2cm}

%%%%%%%%%%%%%%%%%%%%%%%%%%%%%%%%%%%%%%%%%
%%%%%%%%%%%%%%%%%%%%%%%%%%%%%%%%%%%%%%%%%
\subsection{The environment}

% The environment defines which actions the agent can choose and how each action impacts the state of the agent.
% We model the environment as a finite Markov Decision Process.


% The formal definition of a finite Markov decision process goes as follows:

\begin{definition}[Finite Markov decision process]
\label{def: mdp}
A finite Markov decision process consists of
\vspace{-0.3cm}
\begin{enumerate}
    \item a finite set of states $\sset$,
    % with states that the agent can encounter.
    
    \item a finite set of actions $\aset(s)$ for each state $s \in \sset$,
    % with actions that the agent can take in each state.
    
    \item a transition function $\transit : \sset \times \sset \times \aset \to [0,1]$ where
    % \begin{align*}
    %     \transit(s' \mid s, a) = \prob(S_{t+1} = s' \mid S_t = s, A_t = a)
    % \end{align*}
    $\transit(s' \mid s, a)$ denotes the probability of reaching state $s'$ when taking action $a$ in state $s$,
    
    \item a reward function $r : \sset \times \aset \to \real$ where $r(s,a)$ denotes the reward 
    % that the agent receives 
    when taking action $a$ in state $s$,

    \item a discount factor $\discount \in [0,1]$ weighting the importance of immediate versus future rewards.
\end{enumerate}
\end{definition}

% single step
% Thus, when the agent is in a state $s \in \sset$ and takes an action $a \in \aset(s)$, it receives a reward $r(s,a)$ and transitions to a new state $s'$ that is sampled according to the distribution $\transit(s' \mid s, a)$.

% Importantly, the transition from $s$ to $s'$ is independent of any state or action encountered before $s$. This is a fundamental property of MDPs:
% \begin{property}[Markov Property]
%     The future is independent of the past given the present.
% \end{property}
% The Markov Property simplifies the decision-making task, because it guarantees that the impact of a particular action in a given state is independent of how the agent got there.


Decision-making tasks often require many steps, like navigating through a maze or playing chess. The agent then interacts with its environment repeatedly until it enters a terminal state, i.e. a state with no actions.
Such a simulation is called an episode.

\begin{definition}[Episode]
An episode with terminal state $s_T$ is a sequence
% of states, actions and rewards 
\begin{align*}
    (s_t, a_t, r_t)_{0 \leq t < T} \coloneqq (s_0, a_0, r_0, s_1, a_1, r_1, \dots, s_{T-1}, a_{T-1}, r_{T-1})
\end{align*}
such that for every time step $t < T$
\begin{align*}
    a_t \in \aset(s_t) ,
    \quad\quad
    r_t = r(s_t, a_t) ,
    \quad\quad
    \ s_{t+1} \sim \transit(\cdot \mid s_t, a_t) .
\end{align*}
% The terminal state $s_T$ has no actions.
The return of the episode for each state-action pair $(s_t, a_t)$ is
\begin{align*}
    r_t + \discount r_{t+1} + \discount^2 r_{t+2} + \dots + \discount^{T-1-t} r_{T-1} .
\end{align*}
\end{definition}

\vspace{-0.6cm}

%%%%%%%%%%%%%%%%%%%%%%%%%%%%%%%%%%%%%%%%%
% \paragraph{Immediate versus future rewards}

% \begin{remark}[Immediate versus future rewards]
Rewards are discounted at every step of an episode by a factor $\discount$, which balances the importance of immediate versus future rewards.
For instance, when $\discount = 1$, 
immediate and future rewards are equally valued,
    % justifying short-term sacrifices for long-term gains
making it worthwhile to take a low-reward action now to reach states with high rewards later.
In contrast, when $\discount = 0$, only immediate rewards matter.
% \end{remark}

%%%%%%%%%%%%%%%%%%%%%%%%%%%%%%%%%%%%%%%%%
% \paragraph{Stochastic versus deterministic rewards}

%%%%%%%%%%%%%%%%%%%%%%%%%%%%%%%%%%%%%%%%%
% \paragraph{Reward versus cost}

% \mytodo{be more vague, talk more about domains of applications and cite several for each domain}
% This framework covers many real-world tasks, such as
% planning medical treatments~\citep{Schaefer2005ModelingProcesses},
% choosing the temperatures and stirring rates for a bioreactor~\citep{Ungar1991AControl},
% managing a portfolio~\citep{Filos2018ReinforcementManagement},
% multiplying matrices~\citep{Fawzi2022DiscoveringLearning} or playing games~\citep{Schrittwieser2020MasteringModel, Vinyals2019GrandmasterLearning}.


% %%%%%%%%%%%%%%%%%%%%%%%%%%%%%%%%%%%%%%%%%
% \paragraph{Example of a maze}

% For instance, consider the MDP described in \autoref{fig: example MDP}

% represent an MDP
% states are nodes
% actions are edges between nodes

% % \begin{example}[Chess] \label{ex: chess as MDP}
The game of chess is a finite \mdpS.
Namely, the position of all the pieces on the board forms the state of the game. Thus the set $\sset$ contains all the legal board configurations.
% Besides, $\sset$ is finite because the size of the board and the number of pieces is finite.
Further, for each board state $s_t$ a player can choose from a finite set of actions $\aset(s_t)$. As the player plays his move, the board transitions to a new state $s_\tpo$.
Chess is a deterministic \mdpS~since the transition probabilities are either 1 or 0.
\end{example}
% \input{figures/01/example-mdp}

% link up to points discussed above
% - can express as reward or cost
% % - if change discount it is not worth going to (no this is about agent...)

% overall very important but difficult, when given a task to design the right MDP

%%%%%%%%%%%%%%%%%%%%%%%%%%%%%%%%%%%%%%%%%
%%%%%%%%%%%%%%%%%%%%%%%%%%%%%%%%%%%%%%%%%
\subsection{The agent}

The agent is defined by a policy function $\pol : \aset \times \sset \to [0,1]$ where $\pol(a \mid s)$ denotes the probability of taking action $a$ in state $s$.
% A policy $\pol$ thus corresponds to a strategy that the agent follows when interacting with its environment.
If this probability is one for some action in each state, the policy $\pol$ is \textit{deterministic}.
% The policies and deterministic policies make up the function spaces
Such deterministic policies make up the function space
\begin{align}
    \polset \coloneqq 
    \left\{ \pol : \aset \times \sset \to \{ 0,1 \}
    % \mid
    \ \bigg| \
    \forall s \in \sset : \sum_{a \in \aset(s)} \pol(a \mid s) = 1 
    % \tand \pol(a \mid s) \in \{0, 1\} 
    \right\} .
    % \\
    % \polset_D &\coloneqq 
    % \big\{ \pol \in \polset
    % \mid
    % \forall s \in \sset : \pol(a \mid s) \in \{0, 1\}  \big\} .
\end{align}


\begin{remark}[Abuse of notation]
As is common in the literature, we use the notation $\pol(s)$ to refer to the distribution $\pol(\cdot \mid s)$ over $\aset(s)$.
We also insert the distribution $\pol(s)$ in functions with domain $\aset(s)$.
For instance, the expected immediate reward when following policy $\pol$ in state $s$ is
\begin{align*}
    r(s, \pol(s)) \coloneqq \sum_{a \in \aset(s)} \pol(a \mid s) r(s, a) .
\end{align*}
The second ``$s$'' in $r(s, \pol(s))$ is redundant as it is impossible to take actions outside $\aset(s)$ in state $s$. So we shorten the notation to
\begin{align*}
    r(s, \pol) \coloneqq r(s, \pol(s)).
\end{align*}
If the policy $\pol$ is deterministic in state $s$, then $\pol(s)$ directly refers to the action $a$ where $\pol(a \mid s) = 1$.
\end{remark}
% \vspace{-0.3cm}


% Note that the number of policies that are deterministic in all states is finite because the number of state-action pairs is finite. 
% \begin{align}
%     \polset = \big\{ \pol : \forall s \in \sset, \exists a \in \aset(a) : \pol(a \mid s) = 1 \big\} .
% \end{align}

% when we refer to ``the agent'' later, that means the person running the algorithm

% if fix policy on MDP, obtain Markov Chain. read \cite[Chapter~6]{Grimmett2020ProbabilityProcesses} to see if any useful properties

% Another important property is whether a policy is proper or not.
% \begin{definition}[Proper policy]
% A policy is proper if any simulation of the MDP with that policy eventually reaches a terminal state with probability one.
% \end{definition}
% Further, if all the policies in $\polset$ are proper, then the \mdpS~is episodic.
% % \begin{definition}[Episodic MDP]
% % \label{def: episodic mdp}
% % An MDP is episodic if it has at least one terminal state and every policy $\pol \in \polset$ is proper.
% % \end{definition}
% Every simulation of an episodic MDP eventually terminates.


%%%%%%%%%%%%%%%%%%%%%%%%%%%%%%%%%%%%%%%%%
% \paragraph{Performance}

\vspace{-0.6cm}


The agent's performance is quantified by
its expected return in each state.
% when following a policy $\pol$. 
This performance is represented by 
two functions: the \textit{state-value function} and the \textit{action-value function}.
% \vspace{0.3cm}

% \vspace{-0.1cm}

\begin{definition}[Value functions]
The state-value function $v_\pol : \sset \to \real$ measures the expected return
when starting in state $s$ and following policy $\pol$:
\begin{align*} 
% \label{eq: definition state-value function}
    v_\pol(s) 
    \coloneqq 
    \E_{\transit} \big[ r(s_0, \pol) + \discount r(s_1, \pol) + \discount^2 r(s_2, \pol) + \dots \mid s_0 = s \big] .
    % \\
    % &= r(s, \pol(s)) + \discount \sum_{s' \in \sset} \transit(s' \mid s, \pol(s)) \big[ r(s', \pol(s')) + \discount \sum_{s'' \in \sset} \transit(s'' \mid s', \pol(s')) [\dots] \big].
    % \Big( r(s_\tpo, \pol(s_\tpn)) + \discount \sum_{s_\tpt} \transit(s_\tpt \mid s_\tpo, \pol(s_\tpo)) \Big( \dots \Big)\Big) .
    % \nonumber \\
    % 
    % &= \Expect_\pol \Big[R_\tpn + \discount R_\tpo + \dots \mid S_\tpn = s_\tpn \Big] .
\end{align*}
The action-value function $q_\pol : \sset \times \aset \to \real$ measures the expected return when taking action $a$ in state $s$ and then following policy $\pol$:
\begin{align*} 
% \label{eq: definition action-value function}
    &q_\pol(s,a)
    \coloneqq
    \E_{\transit} \big[ r(s_0, a_0) + \discount r(s_1, \pol)
    \\
    &\hspace{4.5cm} + \discount^2 r(s_2, \pol) + \dots \mid s_0 = s, a_0 = a \big] .
    % \\
    % &= r(\satpn) + \discount \sum_{s_\tpo \in \sset} \transit(s_\tpo \mid \satpn) v_\pol(s_\tpo) .
    % \nonumber \\
    % 
    % &= \Expect_\pol\Big[R_\tpn + \discount R_\tpo + \dots \mid S_\tpn = s, A_t = a \Big] .
\end{align*}
\end{definition}
% \vspace{0.2cm}

% \newpage

\vspace{-0.6cm}

For every policy $\pol$ and state $s$, these value functions satisfy
\begin{align}
\label{eq: relate v to q}
    v_{\pol}(s) = q_{\pol}(s, \pol) .
\end{align}

% note that markov property and deterministic policies => recursive definition
% An essential property of the value functions is that they satisfy recursive definitions as for every state it holds that
They also follow a recursive structure. For every state, we see that 
\begin{align}
    \nonumber
    v_\pol(s)
    &\coloneqq \E_{\transit} \big[ r(s_0, \pol) + \discount r(s_1, \pol) + 
    % \discount^2 r(s_2, a_2) + 
    \dots \mid s_0 = s \big]
    \\
    \nonumber
    &= r(s, \pol) + \discount \E_{\transit} \big[ r(s_1, \pol) + 
    \discount r(s_2, \pol) + 
    \dots \mid s_0 = s, a_0 = \pol(s) \big]
    % \\
    % \nonumber
    % &= r(s, \pol) + \discount \sum_{s' \in \sset} \transit(s' \mid s, \pol) \E_\transit \big[ r(s_0, \pol) + \discount r(s_1, \pol) + \dots \mid s_0 = s' \big]
    \\
    \label{eq: bellman eq sval}
    &= r(s, \pol) + \discount \sum_{s_+ \in \sset} \transit(s_+ \mid s, \pol) v_\pol(s_+) ,
\end{align}
and for every state-action pair
\begin{align}
    \nonumber
    q_\pol(s,a)
    &\coloneqq \E_{\transit} \big[ r(s_0, a_0) + \discount r(s_1, \pol) + 
    % \discount^2 r(s_2, a_2) + 
    \dots \mid s_0 = s, a_0 = a \big]
    \\
    \nonumber
    &= r(s, a) + \discount \E_{\transit} \big[ r(s_1, \pol) + 
    \discount r(s_2, \pol) + 
    \dots \mid s_0 = s, a_0 = a \big]
    % \\
    % \nonumber
    % &= r(s, a) + \discount \sum_{s' \in \sset} \transit(s' \mid s, a) \E_\transit \big[ r(s_0, \pol) + \discount r(s_1, \pol) + \dots \mid s_0 = s' \big]
    % \sum_{s' \in \sset} \transit(s' \mid s, a) \bigg[ r(s', \pol(s')) + \discount \sum_{s'' \in \sset} \transit(s'' \mid s', \pol(s')) [\dots] \bigg]
    % \\
    % \nonumber
    % &= r(s, a) + \discount \sum_{s' \in \sset} \transit(s' \mid s, a) v_\pol(s')
    \\
    \label{eq: bellman eq aval}
    &= r(s, a) + \discount \sum_{s_+ \in \sset} \transit(s_+ \mid s, a) q_\pol(s_+, \pol) .
\end{align}
Equations \eqref{eq: bellman eq sval} or \eqref{eq: bellman eq aval}, known as the \textit{Bellman equations}, form a linear system whose solution $v_\pol$ or $q_\pol$ is bounded and unique. This property holds in discounted MDPs where $\discount < 1$,
or in episodic MDPs where any episode under each policy eventually terminates \citep{Bellman2010DynamicProgramming}.
% Since the optimal policy maximises the expected return in every state, the Bellman equations for the optimal policy are:

\vspace{-0.1cm}


Value functions induce a partial order over the set of policies. 
A policy $\pol$ is at least as good as another policy $\pol'$ if its expected return is greater or equal in all states, and thus also in all state-action pairs.
We write $\pol \succeq \pol'$ if and only if $v_{\pi} \ge v_{\pi'}$ or $q_{\pi} \ge q_{\pi'}$.
If, additionally, there exists a state $s$ where $v_{\pi}(s) > v_{\pi'}(s)$, then $\pi$ is better than $\pi'$, denoted $\pi \succ \pi'$.

\vspace{-0.1cm}

% In other words, $\pol \geq \pol'$ if and only if $v_{\pol}(s) \geq v_{\pol'}(s)$ for all states $s \in \sset$ or $q_{\pol}(s,a) \geq q_{\pol'}(s,a)$ for all states $s \in \sset$ and actions $a \in \aset(s)$.
For discounted or episodic MDPs, there always exists at least one policy that is better than or equal to all others.
We denote the set of all such optimal policies by $\mathcal P_*$.
% Furthermore, at least one policy in $\mathcal P_*$ is deterministic 
% and any element in that set by 
By definition, all optimal policies share the same value functions, denoted $v_{\best{\pi}}$ and $q_{\best{\pi}}$.
% 
% defined as $v_{\best{\pol}} = \max_{\pol} v_\pol$ and $q_{\best{\pol}} = \max_{\pol} q_\pol$.
% \begin{align*}
%     % &v_{\best{\pol}}(s) = \max_{\pol} v_\pol(s)
%     &v_{\best{\pol}} = \max_{\pol} v_\pol
%     % \quad \quad
%     &q_{\best{\pol}} = \max_{\pol} q_\pol 
%     % &q_{\best{\pol}}(s,a) = \max_{\pol} q_\pol(s,a) 
% \end{align*}
% % for all states $s \in \sset$ and actions $a \in \aset(s)$.
Likewise, at least one policy is worse than or equal to all others.
% , which we denote by $\worst{\pol}$.
% Its value functions are $v_{\worst{\pol}}$ and $q_{\worst{\pol}}$.

\vspace{-0.1cm}

A policy is optimal if and only if it satisfies the \textit{Bellman optimality equations}. These equations are defined for every state and action as
\begin{align*}
    v_{\best{\pol}}(s)
    % &= \max_{\pol} v_\pol(s)
    % \\
    % &= \max_{\pol} r(s, \pol) + \discount \sum_{s' \in \sset} \transit(s' \mid s, \pol) v_{\pol}(s') 
    % \\
    &= \max_{a \in \aset(s)} r(s, a) + \discount \sum_{s_+ \in \sset} \transit(s_+ \mid s, a) v_{\best{\pol}}(s_+) ,
    \\
    q_{\best{\pol}}(s,a)
    % &= \max_{\pol} q_\pol(s,a)
    % \\
    % &= \max_{\pol} r(s, a) + \discount \sum_{s' \in \sset} \transit(s' \mid s, a) q_{\pol}(s', \pol)
    % \\
    &= r(s, a) + \discount \sum_{s_+ \in \sset} \transit(s_+ \mid s, a) \max_{a_+ \in \aset(s_+)} q_{\best{\pol}}(s_+, a_+) .
\end{align*}
% \begin{definition}[Greedy policy]
\label{def: set of greedy policies}
    A policy $\pol$ is greedy with respect to $q$ if for all states $s$, $q(s, \pol) = \max_a q(s,a)$.
    % \begin{align}
    %     q(s, \pol) = \max_a q(s,a) .
    % \end{align}
\end{definition}

Equation \eqref{eq: relate v to q} implies that if a policy $\pol$ is greedy with respect to its own action values $q_\pol$, meaning that for all states
\begin{align*}
    % \pol(s) = \argmax_{a \in \aset(s)} q_{\pol}(s, a) ,
    q_{\pol}(s, \pol) = \max_{a \in \aset(s)} q_{\pol}(s,a),
\end{align*}
then $\pol$ is optimal since it satisfies the Bellman optimality equation:
\begin{align*}
    v_{\pol}(s) 
    &= q_{\pol}(s, \pol)
    \\
    &= \max_{a \in \aset(s)} q_{\pol}(s, a)
    \\
    &= \max_{a \in \aset(s)} r(s, a) + \discount \sum_{s_+ \in \sset} \transit(s_+ \mid s, a) v_{\pol}(s_+) .
\end{align*}
In other words, every policy in $\mathcal P_*$ is greedy with respect to $\qopt$.
Hence, at least one policy in $\mathcal P_*$ is deterministic, which we refer to as $\pi_*$.

% can also compute the greedy policy using $v_{\best{\pol}}$, but need to know parameters of the MDP such as reward function, transition probabilities and discount factor.

Given that some deterministic policy is optimal and the set of deterministic policies is finite, one could find it by looping through all deterministic policies, solving the Bellman equations and verifying the Bellman optimality conditions.
However, this approach is inefficient as it ignores previously gathered information to guide the search.
Moreover, it is infeasible when the parameters of the MDP are unknown.
% which is often the case in practice. 
In such situations, efficient model-free learning algorithms are necessary.

% \begin{definition}[Advantage function]
% The \textit{advantage function} $a_\pol : \sset \times \aset \to \real$ measures the total discounted reward that an agent expects when choosing an action $a_\tpn$ in a state $s_\tpn$ and then following policy $\pol$, namely
% \begin{align} \label{eq: definition advantage function}
%     a_{\pol}(s, \pol') &= q_{\pol}(s, \pol') - v_{\pol}(s)
%     \\
%     a_{\pol'}(s, \pol) &= q_{\pol'}(s, \pol) - v_{\pol'}(s)
% \end{align}
% \end{definition}
% Interpretation
% For any state $s$, the definition of the state-value functions yields that
% \begin{multline*}
%     v_{\pol}(s) - v_{\pol'}(s) 
%     = r(s, \pol) + \discount \sum_{s' \in \sset} \transit (s' \mid s, \pol) v_{\pol}(s') \\
%     - r(s, \pol') - \discount \sum_{s' \in \sset} \transit (s' \mid s, \pol') v_{\pol'}(s')
% \end{multline*}
% and the advantage functions
% \begin{align*}
%     a_{\pol}(s, \pol') &= r(s, \pol') - r(s, \pol) + \discount \sum_{s' \in \sset} ( \transit (s' \mid s, \pol') - \transit (s' \mid s, \pol) ) v_{\pol}(s')
%     \\
%     a_{\pol'}(s, \pol) &= r(s, \pol) - r(s, \pol') + \discount \sum_{s' \in \sset} ( \transit (s' \mid s, \pol) - \transit (s' \mid s, \pol') ) v_{\pol'}(s')
% \end{align*}
% Note that the advantage is zero for any state where the policies are similar.
% As a result, for any state $s \in \sset$, it holds that
% \begin{align}
% \label{eq: link pol and val}
%     v_{\pol}(s) - v_{\pol'}(s) 
%     &= a_{\pol'}(s, \pol(s)) + \discount \sum_{s' \in \sset} \transit(s' \mid s, \pol) (v_{\pol}(s') - v_{\pol'}(s'))
%     \\
%     &= \discount \sum_{s' \in \sset} \transit(s' \mid s, \pol') (v_{\pol}(s') - v_{\pol'}(s')) - a_{\pol}(s, \pol'(s))
% \end{align}
% These recursive equations indicate that the difference $v_{\pol}(s) - v_{\pol'}(s)$ equals a (possibly infinite) discounted sum of advantage values.

% \begin{theorem}[Banach fixed point theorem] 
\label{thm: banach fixed point}
Let $(\mathcal{X}, \|\cdot\|_\mathcal{X})$ be a metric space, and $\mathsf{F} : \mathcal{X} \to \mathcal{X}$ be a contractive operator. Then,
\begin{itemize}
    \item there is a unique fixed point $x\opt \in \mathcal{X}$ that satisfies $x\opt = \mathsf{F} x\opt$,
    
    \item the sequence $x\new = \mathsf{F} x\old$ converges to $x\opt$ for any $x_0 \in \mathcal{X}$.
\end{itemize}
\end{theorem}

% \begin{remark}
% Since the state-action space is finite, the Bellman operators are equivalent to a linear matrix transformation
% where $r$ and $q$ are vectors with components $r(s,a)$ and $q(s,a)$, respectively, and the matrix $\T_\pol$ is a square matrix that contains the transition probabilities under policy $\pol$ from one state-action pair to another.
% \end{remark}


% \begin{remark}[Reward versus cost]
% Some texts consider a cost function $c$ instead of a reward function $r$. The goal of the agent is then to minimise the total discounted cost it experiences instead of maximising its reward.
% Nevertheless, this modification does not affect the results in this report since a reward-based model can always be turned into a cost-based model by setting $c = -r$, and vice versa.
% \end{remark}

% \textcolor{red}{say assume bounded action values}
% - either discount < 1
% - or markov chain induced by every policy is absorbing
% bc if action values are unbounded, makes little sense 

% \begin{property}
% \label{thm: bellmann of greedy policies}
% If the policies $\pol$ and $\pol'$ are greedy on $q$, then $\B q = \B_{\pol} q = \B_{\pol'} q$.
% \end{property}
% \begin{proof}
% \end{proof}

% $\epsilon$-greedy policy
% greedy policy on $q$
% \begin{align*}
%     \pol(a \mid s) = 
%     \begin{cases}
%         \ \epsilon/|\aset(s)| + 1 - \epsilon  &\text{if} \ a=\argmax_{a'} q(s,a')
%         \\
%         \ \epsilon/|\aset(s)| &\text{otherwise}
%     \end{cases}
% \end{align*}

% An important insight is that MDPs resemble Russian dolls.
% \begin{lemma}
% \label{thm: Russian dolls}
%     For every MDP $\mathcal{M}$ with best and worst policies $\best{\pol}$ and $\worst{\pol}$
%     if
%     it holds that
%     \begin{align}
%         \label{eq: Russian ineq best pol}
%         v_{\best{\pol}'} \leq v_{\best{\pol}}
%         \\
%         v_{\worst{\pol}'} \geq v_{\worst{\pol}}
%         \label{eq: Russian ineq worst pol}
%     \end{align}
%     If $a = \best{\pol}(s)$, then $\worst{\pol}' = \worst{\pol}$ and inequality \eqref{eq: Russian ineq best pol} is strict in at least one state.
%     If $a = \worst{\pol}(s)$, then $\best{\pol}' = \best{\pol}$ and inequality \eqref{eq: Russian ineq worst pol} is strict in at least one state.
% \end{lemma}
% \begin{proof}
% \end{proof}


%%%%%%%%%%%%%%%%%%%%%%%%%%%%%%%%%%%%%%%%%
%%%%%%%%%%%%%%%%%%%%%%%%%%%%%%%%%%%%%%%%%
\vspace{-0.25cm}
\section{Policy iteration}
\vspace{-0.25cm}

Policy iteration (PI) is a value-based, model-free reinforcement learning framework that optimises a policy $\pol$ without knowing the transition or reward functions of the underlying MDP.
Instead, it learns the
impact of each action on future rewards 
% expected return of each action
through trial and error and stores this information in a value function $q$.
PI-based algorithms iteratively refine the policy $\pol$ and the estimate $q$ by alternating between
\textit{policy evaluation}, which updates $q$ to approximate the action values of $\pol$ by simulating the MDP,
and \textit{policy improvement}, which updates $\pol$ to obtain higher returns as anticipated by $q$.
This process enables learning the optimal policy $\best{\pol}$ and its action values $\qopt$ in unknown environments.
This section discusses different approaches to policy evaluation and policy improvement, along with their convergence properties.


% \textit{policy evaluation}, where the estimate $q$ evaluates the action values of current policy $\pol$ by simulating the MDP,
% and \textit{policy improvement}, where the policy $\pol$ is updated to yield higher expected returns as anticipated by $q$.

% Essentially, policy iteration estimates the action values of a policy to infer an order over the actions of each state and, with that, derive a better policy.

% This framework enables effective decision-making in unknown environments.
% Therefore, 
% as its mechanism underpins many reinforcement learning algorithms.
% - do not need to assume anything about the environment
% eg don't know the rules of a game
% Since it many real-world reinforcement learning algorithms use this mechanism
% - MuZero
% - ChatGPT
% much like learning to play a game without knowing the rules.


% policy based or value-based
% optimisation is always alternating between computing something
% There are several types of model-free reinforcement learning algorithms. 
% Value-based algorithms optimise the value function, as a proxy for the policy.
% Policy-based algorithms parameterise and optimise the policy using the gradient of the expected return with respect to the policy parameters.
% Actor-Critic algorithms then combine value-based and policy-based approaches, where the actor learns a policy and the critic learns a value function.
% Some research suggests that value-based and actor critic methods are related \cite{Wagner2013OptimisticResult}


%%%%%%%%%%%%%%%%%%%%%%%%%%%%%%%%%%%%%%%%%
%%%%%%%%%%%%%%%%%%%%%%%%%%%%%%%%%%%%%%%%%
\subsection{Policy evaluation}

Policy evaluation involves estimating the value function of a given policy $\pol$.
Model-free algorithms do this iteratively by simulating the MDP over several steps while following policy $\pol$. After each simulation, they compute an approximation $\Tilde{q}_\pol$ based on the observed rewards and update an estimate $q$ using a stochastic approximation scheme:
\begin{align}
    q_+ = q + \lrate ( \Tilde{q}_\pol - q )
\end{align}
where $\lrate$ is a decreasing learning rate
and $q$ is a function in $\qset$ defined as
\begin{align}
    \qset \coloneqq \big\{ \sset \times \aset \to \real \big\} .
\end{align}
% Under mild assumptions, the estimate $q$ converges to the true action values $q_\pol$ with sufficient simulations.
We typically classify policy evaluation methods based on the length of each simulation, distinguishing between the Monte Carlo method and multi-step temporal difference methods.

% There are two types of policy evaluation methods, each running distinct simulations and producing different estimates for $\Tilde{q}_\pol$.

% estimate the value function of a policy to infer the structure of the MDP bc you compute how a reward is cumulated thgrough by the MDP

% depending on whether the algorithm samples entire or partial episodes.

% This allows the agent to learn and refine value estimates based on observed rewards without requiring a full model of the environment.


%%%%%%%%%%%%%%%%%%%%%%%%%%%%%%%%%%%%%%%%%
\subsubsection{Monte Carlo method}

The Monte Carlo (MC) method simulates entire episodes $(s_t, a_t, r_t)_{0 \leq t < T}$. The state-action pairs visited during the episode are then updated using
\begin{align}
    \Tilde{q}_\pol(s_t, a_t) = r_t + \discount r_{t+1} + \discount^2 r_{t+2} + \dots + \discount^{T-1-t} r_{T-1} .
\end{align}
% If a state-action pair is visited multiple times during an episode, 
% \textit{first-visit MC methods} only update the estimate after the first visit of that state-action pair.
% \textit{Every-visit MC methods}, on the other hand, update the estimate after every visit of that state-action pair.
% \textit{Initial-visit MC methods} only update the first state-action pair of the episode.
% 
% Note that $\E [\Tilde{q}_\pol] = q_\pol$
This method yields an unbiased approximation of the action values
% \begin{align*}
%     \E [\Tilde{q}_\pol] = q_\pol
% \end{align*}
because $q_\pol(s_t,a_t)$ is by definition the expected total discounted reward when taking action $a_t$ in state $s_t$ and then following policy $\pol$.
% Although, it tends to have a large variance due to the randomness of the transitions in the episode.
% As a result, if each state-action pair is visited infinitely often, MC methods converge to $q_\pol$.

% use eligibility traces to avoid having to wait for the end of the episode before doing the update (check how this affects the proofs of convergence) -> transform forward looking method into backwarda looking method


%%%%%%%%%%%%%%%%%%%%%%%%%%%%%%%%%%%%%%%%%
\subsubsection{Temporal difference methods}

Temporal difference (TD) methods simulate $n$ steps $(s_t, a_t, r_t)_{0 \leq t < n}$ and 
% , i.e. $(s_0, a_0, r_0, \dots, s_t)$.
update the state-action pairs visited in the partial episode using the approximation
\begin{align}
\label{eq: nstep approx}
    \Tilde{q}_\pol(s_t, a_t) = r_t + \discount r_{t+1} + \dots + \discount^{n-1-t} r_{n-1} + \discount^{n-t} q(s_{n}, \pol) .
\end{align}
% TD methods are said to ``bootstrap'' because they update the estimate in a state-action pair based on another estimate.
This approximation of $q_\pol$ is biased as it bootstraps on the estimate $q$.
% we say that it bootstraps $q$.
% Note that 
% \begin{align*}
%     \E [\Tilde{q}_\pol] = \B_\pol^n q .
% \end{align*}
% The approximation $\Tilde{q}_\pol$ is thus a biased estimate of $q_\pol$ but, given the contraction of the Bellman operator $\B_\pol$, $\Tilde{q}_\pol$ is, on average, closer to $q_\pol$ than $q$.


% MC and TD methods are related. 
The MC method can be regarded as a TD method with $n = \infty$.
More generally, certain methods use a weighted combination of several multi-step approximations.
For instance, TD($\lambda$) weights each $n$-step temporal difference by $(1-\lambda) \lambda^{n-1}$. So TD(0) corresponds to the 1-step TD method while TD(1) recovers the MC method.
As $\lambda$ increases, longer simulations are assigned more weight. The bias of the TD($\lambda$) approximation thus decreases, but the variance increases as more randomness accumulates from stochastic transitions.
Nevertheless, if every state-action pair is updated infinitely often, all methods converge to the action values $q_\pol$ \citep{Sutton1988LearningDifferences, Dayan1992The}.


% look at MC with eligibility traces instead of just IV/FV/EV
% \textbf{eligibility trace}: temporary record of the occurence of an event, the trace tracks which parameters are eligible to change by learning. This is relates to the credit assignment problem \cite{Singh1996ReinforcementTraces}

% learn from incomplete episodes of experience by bootstrapping,
% update value based on current estimate of value in some successor state (exploits Markov property)
% +ve can learn after every step in episode, from incomplete episodes (can do on-line learning = update during episode)
% +ve works for non-terminating env
% +ve TD target has low variance since only depends on one random actions/transition/rewards
% -ve TD target is biased
% -ve sensitive to initial value function

% Why use MC instead of TD => look at (van Seijen 2016)
% +ve of n-step bootstrapping methods
% + and link to error reduction property (lower bias of target but higher variance)
% -ve of n-step bootstrapping methods
% - require more memory to store n-steps
% - can only update after n-steps
% - require more computations per time steps than 1-step methods


%%%%%%%%%%%%%%%%%%%%%%%%%%%%%%%%%%%%%%%%%
%%%%%%%%%%%%%%%%%%%%%%%%%%%%%%%%%%%%%%%%%
\subsection{Policy improvement}

Policy improvement derives a better policy $\pol_+$ from the action values $q_\pol$. Model-free algorithms achieve this as a consequence of the following theorem.

\label{thm: policy improvement theorem}\begin{theorem}[Policy Improvement Theorem 
% \citep{Sutton2018ReinforcementIntroduction}
% \citep{Howard1960DynamicProcesses}
]
For all policies $\pol$ and $\pol'$, we have $\pol \preceq \pol'$
if for every state $s \in \sset$
\begin{align}
\label{eq: PIT if}
    q_\pol(s, \pol) \leq q_\pol(s, \pol') .
\end{align}
% Then the policy $\pol'$ is at least as good as $\pol$. In other words, 
% % it must obtain a greater or equal expected return from all states 
% % for every state $s \in \sset$ 
% it holds that
% \begin{align}
% \label{eq: PIT then}
%     v_{\pol} \leq v_{\pol'} 
%     \quad \tand \quad
%     q_{\pol} \leq q_{\pol'} .
% \end{align}
% Moreover, if in some state inequality \eqref{eq: PIT if} is strict, then inequality \eqref{eq: PIT then} is strict in that state as well.

% \begin{enumerate}[(a)]
    
%     \item 
%     % $\forall\; s \in \sset \quad q_\pol(s, \pol(s)) \leq q_\pol(s, \pol'(s)) \quad\implies\quad v_\pol \leq v_{\pol'}$
%     $q_\pol(\pol) \leq q_\pol(\pol') \quad\implies\quad v_\pol \leq v_{\pol'}$
    
%     \item 
%     % $\forall\; s \in \sset \quad q_\pol(s, \pol(s)) \geq q_\pol(s, \pol'(s)) \quad\implies\quad v_\pol \geq v_{\pol'}$
%     $q_\pol(\pol) \geq q_\pol(\pol') \quad\implies\quad v_\pol \geq v_{\pol'}$
    
% \end{enumerate}
\end{theorem}
% \begin{proof}
% Equation \eqref{eq: PIT if} implies for every state-action pair $(s,a)$ that
% \begin{align*}
%     \big[ \B_\pol q_\pol \big] (s,a) 
%     &= r(s,a) + \discount \sum_{s' \in \sset} \transit(s' \mid s,a) q_\pol(s', \pol(s')) 
%     \\
%     &\leq r(s,a) + \discount \sum_{s' \in \sset} \transit(s' \mid s,a) q_\pol(s', \pol'(s')) 
%     \\
%     &= \big[ \B_{\pol'} q_\pol \big](s,a) .
% \end{align*}
% \begin{align*}
%     q_\pol (s,a) 
%     &= \E_\transit \big[ r(S_0, A_0) + \discount q_\pol(S_1 \mid \pol) \mid S_0=s, A_0=a \big]
%     \\
%     &\leq \E_\transit \big[ r(S_0, A_0) + \discount q_\pol(S_1 \mid \pol') \mid S_0=s, A_0=a \big]
%     \\
%     &= \E_\transit \big[ r(S_0, A_0) + \discount r(S_1 \mid \pol') + \discount^2 q_\pol(S_2 \mid \pol) \mid S_0=s, A_0=a \big]
%     \\
%     &\leq \E_\transit \big[ r(S_0, A_0) + \discount r(S_1 \mid \pol') + \discount^2 q_\pol(S_2 \mid \pol') \mid S_0=s, A_0=a \big]
%     \\
%     &\dots
%     \\
%     &\leq \E_\transit \big[ r(S_0, A_0) + \discount r(S_1 \mid \pol') + \discount^2 r(S_2 \mid \pol') + \dots \mid S_0=s, A_0=a \big]
%     \\
%     &= q_{\pol'}(s,a) .
% \end{align*}
% Consequently, the monotonicity and the contraction of the Bellman operator $\B_{\pol'}$ indicate that
% \begin{align}
%     \B_\pol q_\pol = q_\pol \leq \B_{\pol'} q_\pol \leq \B_{\pol'}^2 q_\pol \leq \dots \leq \B_{\pol'}^\infty q_\pol = q_{\pol'} 
% \end{align}
% and
% \begin{align}
%     v_\pol
% \end{align}

% % % (Show $\implies$) 
% % (a)
% % Since $q_\pol(s, \pol(s)) \leq q_\pol(s, \pol'(s))$ for all states $s$, we obtain for every state-action pair that
% % \begin{align*}
% %     q_\pol(\sa) 
% %     &= \reward(\sa) + \discount \sum_{s'} \transprob(s' \mid \sa) q_\pol(s', \pol(s')) , \\
% %     &\leq \reward(\sa) + \discount \sum_{s'} \transprob(s' \mid \sa) q_\pol(s', \pol'(s')) = \big[\B_{\pol'} q_\pol \big](\sa) .
% % \end{align*}
% % % $q_\pol(s, \pol(s)) \leq q_\pol(s, \pol'(s))$ is equivalent with $\B_\pol q_\pol \leq $
% % Then, the monotone and contractive properties of the operator $\B_{\pol'}$ imply that
% % \begin{align*}
% %     q_\pol \leq \B_{\pol'} q_\pol \leq \B_{\pol'}^2 q_\pol \leq \dots \leq \lim_{n \to \infty} \B_{\pol'}^n q_\pol = q_{\pol'} .
% %     % v_\pol(s) &= q_\pol(s, \pol(s)) \\
% %     % &\leq q_\pol(s, \pol'(s)) \\
% %     % % &= \mathbb{E}\Big[ R_{t+1} + \discount v_\pol(S_{t+1}) \mid S_t = s, A_t = \pol'(s) \Big] \\
% %     % &= \mathbb{E}_{\pol'}\Big[ R_{t+1} + \discount v_\pol(S_{t+1}) \mid S_t = s \Big] \\
% %     % &\leq \mathbb{E}_{\pol'}\Big[ R_{t+1} + \discount q_\pol(S_{t+1}, \pol'(S_{t+1})) \mid S_t = s \Big] \\
% %     % % &= \mathbb{E}_{\pol'}\Big[ R_{t+1} + \discount \mathbb{E}_{\pol'}\Big[ R_{t+2} + \discount v_\pol(S_{t+2}) \mid S_{t+1} \Big] \mid S_t = s \Big] \\
% %     % &= \mathbb{E}_{\pol'}\Big[ R_{t+1} + \discount R_{t+2} + \discount^2 v_\pol(S_{t+2}) \mid S_t = s \Big] \\
% %     % &\;\;\vdots \\
% %     % &\leq \mathbb{E}_{\pol'}\Big[ R_{t+1} + \discount R_{t+2} + \discount^2 R_{t+3} + \discount^3 R_{t+4} + \dots \mid S_t = s \Big] = v_{\pol'}(s)
% % \end{align*}
% % Thus, for all states it holds that
% % \begin{align*}
% %     v_\pol(s) = q_\pol(s, \pol(s)) \leq q_\pol(s, \pol'(s)) \leq q_{\pol'}(s, \pol'(s)) = v_{\pol'}(s).
% % \end{align*}

% % % (Show $\impliedby$) 
% % % Suppose $q_\pol(s, \pol(s)) \geq q_\pol(s, \pol'(s))$, then
% % % \begin{align*}
% % %     v_\pol(s) &= q_\pol(s, \pol(s)) \\
% % %     &\geq q_\pol(s, \pol'(s)) \\
% % %     % &= \mathbb{E}\Big[ R_{t+1} + \discount v_\pol(S_{t+1}) \mid S_t = s, A_t = \pol'(s) \Big] \\
% % %     &= \mathbb{E}_{\pol'}\Big[ R_{t+1} + \discount v_\pol(S_{t+1}) \mid S_t = s \Big] \\
% % %     &\geq \mathbb{E}_{\pol'}\Big[ R_{t+1} + \discount q_\pol(S_{t+1}, \pol'(S_{t+1})) \mid S_t = s \Big] \\
% % %     % &= \mathbb{E}_{\pol'}\Big[ R_{t+1} + \discount \mathbb{E}_{\pol'}\Big[ R_{t+2} + \discount v_\pol(S_{t+2}) \mid S_{t+1} \Big] \mid S_t = s \Big] \\
% % %     &= \mathbb{E}_{\pol'}\Big[ R_{t+1} + \discount R_{t+2} + \discount^2 v_\pol(S_{t+2}) \mid S_t = s \Big] \\
% % %     &\;\;\vdots \\
% % %     &\geq \mathbb{E}_{\pol'}\Big[ R_{t+1} + \discount R_{t+2} + \discount^2 R_{t+3} + \discount^3 R_{t+4} + \dots \mid S_t = s \Big] = v_{\pol'}(s)
% % % \end{align*}

% % (b) As in part (a), we obtain that
% % \begin{align*}
% %     q_\pol \geq \B_{\pol'} q_\pol \geq \B_{\pol'}^2 q_\pol \geq \dots \geq \lim_{n \to \infty} \B_{\pol'}^n q_\pol = q_{\pol'} .
% % \end{align*}
% % Thus, for all states it holds that
% % \begin{align*}
% %     v_\pol(s) = q_\pol(s, \pol(s)) \geq q_\pol(s, \pol'(s)) \geq q_{\pol'}(s, \pol'(s)) = v_{\pol'}(s).
% % \end{align*}

% \end{proof}

\vspace{-0.6cm}

For discounted and episodic MDPs, the Policy Improvement Theorem guarantees that every policy $\pi'$ that satisfies condition \eqref{eq: PIT if} is at least as good as policy $\pol$.
However, model-free algorithms only approximate $q_\pol$ in the limit rather than computing it exactly.
Therefore, a trade-off emerges depending on the number of policy evaluation steps that precede an improvement step.


%%%%%%%%%%%%%%%%%%%%%%%%%%%%%%%%%%%%%%%%%
\subsubsection{The non-optimistic method}

The non-optimistic method updates the policy once the estimate $q$ has converged to $q_\pol$.
Consequently, it can compute a new policy $\pol_+$ that satisfies condition \eqref{eq: PIT if} exactly and is therefore at least as good as $\pol$.
However, waiting for complete convergence is computationally expensive and unnecessary because we can compute the same policy $\pol_+$ provided the estimate $q$ is ``sufficiently close'' to $q_\pol$.

%%%%%%%%%%%%%%%%%%%%%%%%%%%%%%%%%%%%%%%%%
\subsubsection{Optimistic methods}

\vspace{-0.2cm}

Optimistic methods update the policy after a finite number of evaluation steps.
The policy $\pol_+$ is computed such that $q(s, \pol) \leq q(s, \pol_+)$. It is not necessarily better than $\pol$ as $q$ may not have converged to $q_\pol$ yet,
but sacrificing this certainty can speed up the overall optimisation process.

% Regardless of the method used, we require in the limit that the policy $\pol$ is greedy with respect to the estimate $q$.
% This ensures that the algorithm stabilises once it reaches the optimal policy $\best{\pol}$ and action values $\qopt$.


\vspace{-0.1cm}

%%%%%%%%%%%%%%%%%%%%%%%%%%%%%%%%%%%%%%%%%
%%%%%%%%%%%%%%%%%%%%%%%%%%%%%%%%%%%%%%%%%
\subsection{Greedy in the limit with infinite exploration}
\label{sec: glie}

\vspace{-0.2cm}


The degrees of bootstrapping during evaluation steps and optimism in the improvement steps 
define two dimensions along which all PI-based algorithms lie.
The four corners spanning this space are 1-step TD optimistic and non-optimistic methods and MC optimistic and non-optimistic methods.

The general framework of all PI-based algorithms is the same, but implementations vary.
For instance, the estimate $q$ can be updated offline at the end of each episode as in \citep[p. 99-101]{Sutton2018ReinforcementIntroduction},
or online using eligibility traces \citep{Singh1996ReinforcementTraces}.
The approximation $\Tilde{q}_\pol$ can be computed on-policy as in Sarsa($\lambda$) \citep{Sutton1996GeneralizationCoding},
or off-policy as in \citep[p. 111]{Sutton2018ReinforcementIntroduction}, Retrace($\lambda$) \citep{Munos2016SafeLearning} and various $Q(\lambda)$ algorithms \citep{Harutyunyan2016QCorrections, Sutton2014AEquivalence, Peng1996IncrementalQ-Learning, Watkins1989LearningRewards}.
The approximation can also be averaged over several simulations to filter out stochastic perturbations as in $Q(\sigma)$ \citep{DeAsis2018Multi-StepAlgorithm}.

Regardless of their implementation, all algorithms must follow two essential principles formally known as GLIE.
\vspace{-0.1cm}

\begin{definition}[Greedy in the limit with infinite exploration]
\label{def: glie}
Greedy in the limit with infinite exploration (GLIE) implies that
\vspace{-0.4cm}
\begin{enumerate}
    \item all state-action pairs are updated infinitely many times;

    \item the policy $\pol$ converges to a greedy policy with respect to $q$.
\end{enumerate}
\end{definition}
\vspace{-0.5cm}
The first condition ensures sufficient exploration of the state-action space, while the second guarantees that the optimal policy $\best{\pol}$ and its action values $\qopt$ form a stable fixed point for the algorithm.


% these algorithms are everywhere and have derivatives
% - value iteration and its stochastic version q-learning
% - TD(lambda) interpolates between TD and MC learning => TD(0) = VI, TD(1) offline update = MCES
% - Monte Carlo Exploring starts
% - monte carlo tree search
% - Sarsa
% - AlphaGo, AlphaZero, MuZero
% look at these algorithms in expectation (dynamic programming) or sampled version (this article)


% \begin{table}[h]
% \centering
% \begin{tabular}{|c|c|c|}
% \hline
% % \multirow{2}{*}{Amount of Bootstrapping} & \multicolumn{2}{c|}{Amount of Optimism} \\ \cline{2-3} 
%  & $t < \infty$ & $t=\infty$ \\ \hline
% $n < \infty$ & TD optimistic policy iteration & TD policy iteration \\ \hline
% $n = \infty$ & MC optimistic policy iteration & MC policy iteration \\ \hline
% \end{tabular}
% \caption{Value-based model-free reinforcement learning}
% \label{tab: rl-algorithms}
% \end{table}

%%%%%%%%%%%%%%%%%%%%%%%%%%%%%%%%%%%%%%%%%
%%%%%%%%%%%%%%%%%%%%%%%%%%%%%%%%%%%%%%%%%
\subsection{Convergence}

The convergence of PI-based algorithms
% have been extensively studied.
% as its mechanism underpins many reinforcement learning algorithms. These properties 
% They 
depends primarily on the level of optimism during policy improvement and
secondarily on whether all state-action pairs are updated after each episode using noiseless approximations,
or only a subset of pairs using noisy approximations.

% synchronously using noiseless approximations
% or asynchronously using noisy ones.

% the degree of bootstrapping during policy evaluation.

%%%%%%%%%%%%%%%%%%%%%%%%%%%%%%%%%%%%%%%%%
\subsubsection{Non-optimistic methods converge}

Non-optimistic methods always converge to an optimal policy and its action values under GLIE. This directly follows from the Policy Improvement Theorem that guarantees that non-optimistic methods generate a sequence of monotonically improving policies.

%%%%%%%%%%%%%%%%%%%%%%%%%%%%%%%%%%%%%%%%%
\subsubsection{Optimistic methods converge under synchronous updates}

% The convergence of optimistic methods is more complex.

\cite{Rothblum1979IteratedProcesses} and later \cite{Canbolat2013ApproximateProcesses} showed that optimistic methods converge to the optimal solution if all state-action pairs are updated after each episode using noiseless approximations and a constant learning rate $\lrate=1$.
However, these conditions require either access to the exact parameters of the MDP or averaging over many simulations, making them impractical in most real-world scenarios.


%%%%%%%%%%%%%%%%%%%%%%%%%%%%%%%%%%%%%%%%%
\subsubsection{Certain optimistic methods converge under asynchronous updates}

When state-action pairs are updated asynchronously using noisy simulations and decreasing learning rates, several counterexamples have demonstrated that optimistic methods can converge to suboptimal fixed points \citep{Williams1993AnalysisSystems, Bertsekas2010Williams-BairdIteration, Bertsekas1996Neuro-DynamicProgramming, Wang2022OnLearning}.


Despite these challenges, certain optimistic variants retain strong convergence properties.
For instance, algorithms that use 1-step TD policy evaluation such as $Q$-learning, Sarsa and Expected Sarsa converge to the optimal solution under the GLIE assumptions \citep{Watkins1992Q-learning, Jaakkola1993ConvergenceAlgorithms, Singh2000ConvergenceAlgorithms, vanSeijen2009ASarsa}.
These results follow from a contraction argument that demonstrates resilience to asynchronous noisy updates.


% The convergence of such as is more complex.
The behaviour of multi-step TD methods is more complex because the contraction argument does not directly generalise. Moreover, in the limit where the policy $\pol$ is greedy on the estimate $q$, the approximation $\Tilde{q}_{\pol}$ computed in \eqref{eq: nstep approx} becomes discontinuous in $q$.
Nevertheless, some optimistic methods that use TD($\lambda$) or MC approximations
converge to the optimal solution under specific conditions,
such as when all state-action pairs are updated equally frequently \citep{Tsitsiklis2002OnIteration, Liu2021OnStarts, Chen2018OnProblem},
when the MDP exhibits a feedforward structure \citep{Wang2022OnLearning, Lubars2021OptimisticStructure},
or when the weighting factor $\lambda$ is sufficiently small compared to the discount factor $\discount$ \citep{Harutyunyan2016QCorrections}.
More complex multi-step algorithms manage to recover contraction arguments 
by ingeniously combining several multi-step approximations \citep{Munos2016SafeLearning},
or using a look-ahead mechanism that offsets the undesirable effects of optimistic improvement \citep{Winnicki2023OnEvaluation}.
% \citet[p. 99]{Sutton2018ReinforcementIntroduction} conjecture that 


% - asynchronous updates with particular initial conditions \cite{Bertsekas2010DistributedProgramming}
% - linking it to other frameworks and using their convergence results \cite{Wagner2013OptimisticResult}

% Furthermore, there is 
% also whole literature on the effect of approximating,
% but we restrict ourselves to the tabular case. In other words, assume that we can represent and store the estimate $q$ exactly

%%%%%%%%%%%%%%%%%%%%%%%%%%%%%%%%%%%%%%%%%
\subsubsection{Convergence of MC-O-PI remains an open problem}

Three of the four variants that delimit the space of PI-based algorithms converge to the optimal solution under general conditions, namely the 1-step TD and MC non-optimistic variants, and the 1-step TD optimistic variant.
However, the MC optimistic variant (MC-O-PI) may converge to suboptimal solutions even in its simplest form \citep[Example~5.12]{Bertsekas1996Neuro-DynamicProgramming},
raising the question: \textit{why} and \textit{when} does convergence to optimality break down for multi-step optimistic policy iteration algorithms?

\citet[p.~13]{Sutton1999OpenLearning} emphasises the importance of understanding the convergence of MC-O-PI, saying that
``[i]t is hard to imagine any RL method simpler or more likely to converge than [MC-O-PI] ...
% yet there remain no proof of asymptotic convergence to $\qopt$. 
While this simplest case remains open we are unlikely to make progress on any [multi-step TD-O-PI algorithm].''
\citet[p.~99]{Sutton2018ReinforcementIntroduction} later reinforce this by labelling the problem as
``one of the most fundamental open theoretical questions in reinforcement learning.''


% \textcolor{red}{integrate}
% Furthermore, understanding tabular MC-O-PI, then helps with understanding in general opi
% which is necessary to understand with approximation
% which is linked to policy gradient methds \citep{Wagner2013OptimisticResult}
% which are the learning methods used to train real world agenst like large language models and robotics.



\newpage
Understanding the asymptotic behaviour of MC-O-PI would thus 
yield significant theoretical and practical results.
% It addresses open questions in the field and helps establish implementation guidelines for other PI variants.
% despite its apparent simplicity: continuously updating an estimate $q$ towards the action values of the policy $\pol$ greedy on $q$.
Therefore, we study MC-O-PI with asynchronous noisy updates, aiming to determine necessary and sufficient conditions on its parameters that guarantee convergence to optimality.
% independently of the structure of the MDP.
Importantly, we do not seek conditions on the MDP parameters as these cannot be verified when the model of the environment is unknown.
% We intend to address fundamental open questions in the field 
% such as the conjecture of \citet[p. 99]{Sutton2018ReinforcementIntroduction}
% and establish implementation guidelines for the algorithm.


% related algorithms, such as Monte Carlo exploring starts in \citep[Chapter~5]{Sutton2018ReinforcementIntroduction}
% and GLIE Monte Carlo control \citep{Silver2015LecturesLearning}.


% The convergence analysis of Reinforcement Learning algorithms relies on the action values $q_k$ and the policy $\pol_k$.
% For instance, the convergence proof of Value Iteration relies on the contraction of $q_k$ towards $q_{\best{\pol}}$,
% while the convergence proof of Policy Iteration shows that the sequence $\pol_k$ monotonically improves.
% Since MC-OPI tracks both $q_k$ and $\pol_k$, there is the possibility to study the convergence in the space of action values or the space of policies.


% The reason why convergence of mces is more diffficult to study is bc due to optimism and no bootstrapping new phenomena appear on the boudary, phenomena which are not present when do PI (no optimism so jump over boundary), or VI (100percent bootstrapped so on boundary r+Bq is always the same and cannot have fixed points)

% \vspace{-0.3cm}

%%%%%%%%%%%%%%%%%%%%%%%%%%%%%%%%%%%%%%%%%
%%%%%%%%%%%%%%%%%%%%%%%%%%%%%%%%%%%%%%%%%
\section{Contributions}

% \vspace{-0.3cm}

% why it's relevant
% why it's novel

% exact mcopi
Convergence proofs for PI-based algorithms 
% model them as stochastic difference equations and 
typically rely on one of two techniques.
They either show that the policy improves monotonically using the Policy Improvement Theorem,
% as for non-optimistic methods,
or prove that the estimate approaches the optimal action values using a contraction argument.
% as for 1-step TD optimistic methods.
Unfortunately these techniques do not generally extend to MC-O-PI.
% multi-step TD optimistic methods. 
Therefore previous research imposed conditions either on
the MDP parameters to recover monotone policy improvement,
or
on the MC-O-PI parameters to obtain a contraction argument.
% and show convergence of $q$ starting from the last actions in the decision tree and progressing backwards.


Instead of analysing the noisy MC-O-PI updates directly,
we study the underlying mean-field dynamics.
By averaging out noise, we can track the evolution of the value estimate $q$ and the policy $\pi$. The latter was previously infeasible because the impact of noise on the evolution of the policy is unclear.


We find that MC-O-PI's asymptotic behaviour primarily depends on the probability of updating each state-action pair after an episode.
This probability is, first, determined by the episode's initial state-action pair.
% that defines which state-action pairs of the episode are updated.
% 
A common assumption on the initial state-action pairs is \textit{exploring starts}, which ensures that every state-action pair has a non-zero chance of being the episode's starting point.
Second, the update probability depends on the update function.
The two main update functions are \textit{initial-visit}, which updates only the first state-action pair of the episode,
and \textit{first-visit}, which updates all visited pairs on their first appearance in the episode.
% there is also and Every Visit variant, but we do not consider this.
By analysing the mean-field dynamics for different configurations of update probabilities, we derive the following insights.


% %%%%%%%%%%%%%%%%%%%%%%%%%%%%%%%%%%%%%%%%%
% \subsubsection{Propose an alternative to stochastic approximation techniques}


% Instead of analysing the noisy MC-O-PI updates,
% % using traditional stochastic approximation techniques, 
% we study the deterministic system that emerges when averaging out noise.
% We show that if the optimal action values are globally asymptotically stable for this system, MC-O-PI converges to $\qopt$ with probability one. 

% This approach allows us to establish convergence from either the evolution of the estimate $q$ or the policy $\pol$.
% The latter was previously infeasible because the impact of noise on the evolution of the policy is unclear.
% By combining insights from both sequences, we obtain a more thorough understanding of the algorithm, which allows us to extend all existing convergence proofs while providing more intuitive explanations.


% This approach also shows that MC-O-PI's behaviour mainly depends on the probability of updating each state-action pair after an episode, which is determined by the episode's initial state-action pair and the update function.
% % that defines which state-action pairs of the episode are updated.
% % 
% A common assumption on the initial state-action pairs is \textit{exploring starts}, which ensures that every state-action pair has a non-zero chance of being the episode's starting point.
% The two main update functions are \textit{initial-visit}, which updates only the first state-action pair of the episode,
% and \textit{first-visit}, which updates all visited pairs on their first appearance in the episode.
% % there is also and Every Visit variant, but we do not consider this. 

% Overall, we study the evolution of the estimate $q$ and the policy $\pol$ generated by the deterministic system that underlies MC-O-PI 
% for different configurations of update probabilities.

% % In the process, we show that the robustness theory developed by \cite{Kellett2004SmoothInclusions} for bounded perturbations also applies to random noise with a bounded variance and decreasing learning rate that satisfies the Robbins-Monro conditions.


%%%%%%%%%%%%%%%%%%%%%%%%%%%%%%%%%%%%%%%%%
\subsubsection{Extend convergence proofs under uniform updates}

% we find the underlying mean field generates monotonically improving policies
% however cannot use ODE method or combined stability-ODE method to directly generalise convergnce to stochasti cprocess,
% therefore we build an argument from first principles, based on teh insights of the mean field


% uniform bias
MC-O-PI is known to converge under the strong assumption that all state-action pairs are updated with the same frequency on average,
as is the case when combining a uniform sampling distribution over all state-action pairs with initial-visit updates \citep{Tsitsiklis2002OnIteration, Chen2018OnProblem, Liu2021OnStarts}.
However, this uniformity is unrealistic as the full set of states is typically unknown in large problems.


We show that convergence still holds under the milder condition that update probabilities are uniform over just the actions in each state.
% because the sequence of policies improves monotonically when noise is averaged out.
In other words, with initial-visit updates, the initial state of each episode can be sampled according to any distribution as long as the initial action in that state is chosen uniformly.
This algorithm is implementable when each state offers a reasonable number of actions.
% more practical than sampling uniformly over the entire state-action space as the action space of each state is typically smaller than the state space.

This result suggests that the exploring starts arising from uniform sampling over all state-action pairs are not \textit{necessary} for convergence.
Instead, learning can prioritise different states at different times by sampling the initial state of each episode accordingly.


To obtain this result, we extend the combined stability-ODE method to accommodate the discontinuous dynamics of MC-O-PI \citep[Section~5.4]{Kushner2003StochasticApplications}.
This opens up new possibilities, beyond contraction-based analysis, to study the convergence of optimistic policy iteration algorithms.





% This algorithm is implementable when there are a small number of actions available in each state, 
% and has strong convergence properties, proceeding via a sequence of improving policies in the limit as noise is averaged out.
% link to IV MCES 
% Furthermore, this result highlights that exploring starts is not \textit{necessary} to guarantee convergence, as long as all actions of each state have the same update probability.


%%%%%%%%%%%%%%%%%%%%%%%%%%%%%%%%%%%%%%%%%
\subsubsection{Describe how non-greedy updates prevent convergence}

When updates are non-uniform, MC-O-PI is known to sometimes converge to a suboptimal solution.
These counterexamples all have a strong bias against updating the greedy actions of the estimate $q$.
% at every time step.


We construct a similar example using the simplest possible MDP with one state and two actions.
This example still fails to converge to the optimal solution when sampling initial state-action pairs according to exploring starts, showing that exploring starts is not \textit{sufficient} to guarantee convergence either.

By describing the bifurcation that creates suboptimal fixed points, we demonstrate that a non-greedy bias in one or more states allows solutions to oscillate indefinitely between policies. This endless switching produces a stable suboptimal fixed point as learning rates decrease.
% explain why a non-greedy bias in one or more states often prevents solutions from converging to the optimal solution.
Reversing this analysis, we obtain sufficient conditions such that MC-O-PI converges to the optimal solution for arbitrary update probabilities. These conditions strictly extend the results of \citet{Wang2022OnLearning}.
% and \citet{Lubars2021OptimisticStructure}.


%%%%%%%%%%%%%%%%%%%%%%%%%%%%%%%%%%%%%%%%%
\subsubsection{Conjecture convergence under greedy updates with noisy estimates}

% greedy bias
% The bifurcation analysis highlights that if MC-O-PI updates greedy actions more frequently in all states, solutions cannot oscillate indefinitely between two policies.
% Nevertheless, w

We construct two counterexamples where initial- and first-visit MC-O-PI oscillate indefinitely between suboptimal policies even though greedy actions are updated more frequently in all states.
Yet solutions remain trapped because either the learning rate cancels out the greedy bias, or the absence of noise prevents solutions from visiting the optimal policy.
% The point of convergence in these examples is particularly pernicious as, in the limit, all actions in each state appear equivalent but an unbounded return is available from an alternative policy that is never found.

These examples suggest that if the conjecture of \citet[p. 99]{Sutton2018ReinforcementIntroduction} holds, it is due to the combination of stochastic episode simulations
and first-visit updates rather than either factor alone.
Furthermore, in combination with the non-greedy analysis, this shows that convergence under initial-visit updates can only be guaranteed when update probabilities are uniform.


% a counterexample where IV MC-O-PI oscillates indefinitely between six suboptimal policies, even when the greedy actions are updated more frequently and under exploring starts.

% We also provide a counterexample where FV MC-O-PI fails to converge when noise caused by episode simulations is averaged out. We note, however, that the algorithm converges to the optimal solution with noisy updates.


We further conjecture that if the greedy action in each state is updated more frequently than the other actions in that state, and the learning rates do not counteract this bias, adding centred noise to each episode's returns prevents MC-O-PI from converging to suboptimal fixed points. 
% Such noise can arise naturally in stochastic episode simulations.
% or be added artificially when the MDP is deterministic.
% suboptimal fixed points are not an issue in practice. If the algorithm appears to converge to one, adding centred noise to the MC approximations guarantees solutions eventually leave the suboptimal fixed point.
% prevents solutions from converging to the suboptimal fixed point.
% This noise can naturally arise from stochastic episode simulations.
This noise, which can arise naturally in stochastic episode simulations, ensures that solutions visit the optimal policy infinitely often and the greedy bias prevents them from consistently leaving that policy.
% 
These results allow us, with initial- and first-visit updates, to sample the initial state of each episode according to any distribution as long as the initial action in that state is the current greedy action at least as often as any other action. 
% Under these conditions, MC-O-PI cannot converge to a suboptimal fixed point. 

% It is also more data efficient as it allows updating multiple state-action pairs after each episode instead of only the initial one.

% This situation occurs for instance 
% % when combining uniform update probabilities over all action in each state with first-visit updates,
% for FV MC-O-PI when each episode is simulated using an epsilon-greedy policy, under the GLIE assumptions.



% Overall, our analysis does not depend on how the MC approximation $\Tilde{q}_\pol$ is computed at every step and when the estimate $q$ is updated.
% \mytodo{remove all references and keep it short}
% In other words, the approximation can be computed using either an on-policy approach as in the algorithms on pages 99 and 101 of \citep{Sutton2018ReinforcementIntroduction},
% or an off-policy approach as on page 111.
% Updates can be performed offline at the end of each episode or online using eligibility traces. Sarsa($\lambda$) \cite[Lecture 5, p. 29]{Silver2015LecturesLearning}

% We only require that the approximation is unbiased and has a bounded variance, and that the policy is updated at the end of each episode.
% Furthermore, we assume that the estimate $q$ is represented exactly, with no approximations.



% not exactly fv mces, need to guarantee that bias is greedy for optimal policy
% Both FV MCES and FV GLIE MC are conjectured to always converge to the optimal solution. 
% These algorithms are more efficient but are also more difficult to analyse due to the complications of FV sampling (which also constrains the choice of overall sampling distributions). In both cases, after the sampling of the initial state action pair the updating of subsequent actions will always be strongly biased towards the greedy action. In the case of GLIE MC the initial choice is also be biased towards the greedy action. For these algorithms we provide partial answers (and insight).


% relate to barto and sutton
% That is, the conjecture in Barto and Sutton, as stated, may be true (and we think it likely) but, if so it depends both on a choice made for efficiency (FV) and noise rather than the structure of the problem itself. In our example (perhaps? - this would be nice if true) convergence in fact occurs faster for IV MCES with unform action selection.
% % relate to GLIE
% e.g. eps-greedy sampling of initial state action pair as in FV GLIE MC

%%%%%%%%%%%%%%%%%%%%%%%%%%%%%%%%%%%%%%%%%
%%%%%%%%%%%%%%%%%%%%%%%%%%%%%%%%%%%%%%%%%
\subsubsection{Structure}


% This thesis is structured as follows.
\autoref{ch: mcopi} details MC-O-PI, its properties, existing convergence proofs and counterexamples.
\autoref{ch: exact mcopi} describes the impact of the update probabilities on the mean-field dynamics.
% derives the mean-field system that arises when averaging out noise and shows that if it converges to $\qopt$, so does MC-O-PI.
The next chapters examine how biased update probabilities affect MC-O-PI's asymptotic behaviour.
% through new counterexamples and convergence proofs.
\autoref{ch: unbiased} proves that solutions converge to $\qopt$ when updates are uniform over all actions in each state.
\autoref{ch: non-greedy bias} explains why solutions often converge to suboptimal fixed points when updates are biased towards non-greedy actions in one or more states.
\autoref{ch: greedy bias} argues that solutions converge to $\qopt$ when updates are biased towards greedy actions in each state and the approximations are sufficiently noisy.
% and MC-O-PI relies on stochastic episode simulations.
Finally, \autoref{ch: conclusion} summarises our results and suggests directions for future work.
